\chapter{Einleitung}
\label{ch:einleitung}

\section{Motivation}
% TODO: Flexible Automatisierung, kleine Losgrößen, Bauteile ohne feste Zuführung,
% Qualitätssicherung inline statt stichprobenartig.

\section{Aufgabenstellung}
Ziel des Projekts ist ein Demonstrator, der ein Bauteil mit beliebiger Position und
Orientierung in einem durch vier ArUco-Marker begrenzten Arbeitsbereich erkennt, mit einem
Roboter greift und an einer Prüfstation ablegt. Dort wird eine Seriennummer ausgelesen,
ein Lochdurchmesser vermessen und das Vorhandensein einer Kerbe geprüft. Abhängig vom
Prüfergebnis sortiert der Roboter das Bauteil als Gut- oder Schlechtteil aus.

\section{Aufbau der Arbeit}
\Cref{ch:grundlagen} stellt die theoretischen Grundlagen vor. In \cref{ch:systemkonzept}
wird das Gesamtkonzept beschrieben, die \cref{ch:objekterkennung,ch:robotersteuerung,ch:sichtpruefung,ch:kigreifen}
behandeln die Umsetzung der einzelnen Module. \Cref{ch:evaluation} bewertet das System
anhand von Messreihen, \cref{ch:fazit} fasst die Ergebnisse zusammen.
